\chapter{Epidural Blood Patch}

\section*{Overview}
An epidural blood patch is the injection of the patient's own blood into the epidural space to seal a cerebrospinal fluid (CSF) leak. It is the definitive treatment for post-dural-puncture headache.
\section*{Alternative Names}
Blood patch; autologous epidural blood patch
\section*{Principle}
Autologous blood injected into the epidural space forms a clot that patches the dural puncture site, stopping the continued leak of CSF. Restoring CSF volume and pressure relieves the traction headache that results from intracranial hypotension.
\section*{History}
The epidural blood patch was introduced by Charles B. Gormley in 1960 and rapidly became the standard treatment for post-dural-puncture headache after the report of larger series in the 1970s.
\section*{Why It Is Done}
Ordered for post-dural-puncture headache that is severe, persists despite conservative therapy such as bed rest, hydration, and caffeine, or interferes with daily activity, and that follows lumbar puncture, spinal anesthesia, or epidural catheter placement.
\section*{Is This a Primary or Secondary Test or Procedure}
Primary therapeutic procedure for post-dural-puncture headache.
\section*{How It Is Performed}
Blood is drawn from a peripheral vein under sterile conditions. The patient is positioned prone or lateral, and a needle is placed into the epidural space at the level of or below the original dural puncture, often with fluoroscopic guidance. Fifteen to thirty milliliters of blood are injected slowly and the needle is withdrawn.
\section*{Preparation}
Informed consent and review of bleeding risk and any active infection. The patient should be volume-replete, and blood is collected with aseptic technique before injection.
\section*{Contraindications and Precautions}
Contraindications include local infection at the puncture site, coagulopathy or anticoagulation, patient refusal, and fever or suspected bacteremia. Precaution: placement above the original puncture site may be less effective.
\section*{Risks}
Immediate backache and leg pain at injection, headache recurrence if the patch fails, and rarely infection, bleeding, or nerve injury. Repeat patches may be needed in a minority of cases.
\section*{Aftercare and Recovery}
The patient lies flat for one to two hours and is advised to avoid heavy lifting and straining for several days. Headache typically resolves within hours to a day.
\section*{Results and Interpretation}
Resolution or marked improvement of the headache and ability to resume upright activity indicate success. Persistent headache suggests an ongoing or additional leak, prompting imaging or a repeat patch.
\section*{Expected Results}
Complete resolution of postural headache with return to normal upright posture without symptoms.
\section*{Follow-up}
Patients with recurrent symptoms may undergo MRI to assess for persistent CSF leak and may require a second blood patch or further investigation of intracranial hypotension.
\section*{References}
\begin{enumerate}
  \item MedlinePlus. Epidural blood patch. U.S. National Library of Medicine; 2025.
  \item Current Medical Diagnosis and Treatment. McGraw-Hill; 2025.
\end{enumerate}

\clearpage
